%% Cover Letter — RNAAS Submission
%% Pierre-Olivier Després Asselin
%% To be uploaded alongside the research note during submission

\documentclass[12pt]{article}
\usepackage[margin=2.5cm]{geometry}
\usepackage{hyperref}
\usepackage{parskip}

\begin{document}

\begin{flushright}
  Pierre-Olivier Despr\'{e}s Asselin \\
  Independent Researcher \\
  Montr\'{e}al, Qu\'{e}bec, Canada \\
  \href{mailto:pierreolivierdespres@gmail.com}{pierreolivierdespres@gmail.com} \\
  ORCID: 0009-0009-7611-4550 \\[6pt]
  April 2026
\end{flushright}

\bigskip

\noindent
To the Editors, \\
\textit{Research Notes of the American Astronomical Society}

\bigskip

Dear Editors,

I am submitting for your consideration a Research Note entitled
\textbf{``A Mass-Dependent Transition Radius in Galaxy Rotation
Curves: Evidence for a Universal Scaling Law from the SPARC Sample,''}
for publication in \textit{Research Notes of the AAS}.

\bigskip

\textbf{Purpose of this Research Note.}
This short contribution reports a \emph{purely empirical} finding: a
tight power-law correlation between an optimal transition radius
$r_c$ and the baryonic mass $M_\mathrm{bary}$ across 103 galaxies
from the SPARC catalogue (Lelli, McGaugh \& Schombert 2016). The
relation
\[
    r_c(M_\mathrm{bary}) = 2.6 \times
      (M_\mathrm{bary} / 10^{10}\,M_\odot)^{0.56}\ \mathrm{kpc}
\]
is confirmed with Pearson $r = 0.768$ and $p = 3\times10^{-21}$
($N = 103$). The Note focuses exclusively on this scaling relation and
its potential as a falsifiable test with near-term data (Euclid, DESI,
Rubin Observatory).

\bigskip

\textbf{Scope compliance.}
This submission is intentionally limited to a single observational
result with no theoretical derivation, consistent with the mission of
RNAAS to publish short, observationally-driven results. The broader
theoretical framework that motivated searching for this relation is
described in a separate, longer manuscript submitted to MNRAS, but is
\emph{not} the subject of the present Research Note.

\bigskip

\textbf{Length and format.}
The Note is approximately 1,050 words (target $\le$ 1,500) and
includes one table, consistent with the RNAAS format requirements.

\bigskip

\textbf{Originality.}
The Note has not been published or submitted elsewhere. The empirical
result is the author's own work based on independent fitting of
publicly available SPARC data.

\bigskip

\textbf{Open science.}
All fitting scripts, intermediate per-galaxy results, and the
regression output are available at
\url{https://github.com/chronos717313/Mastery-of-time}
(DOI:~\href{https://doi.org/10.5281/zenodo.18287042}{10.5281/zenodo.18287042}).
Any researcher can reproduce the correlation from the SPARC catalogue
in approximately one hour of analysis time.

\bigskip

\textbf{Financial support and publication charges.}
I am an independent researcher without institutional or grant funding.
I respectfully request consideration for the RNAAS Publication Support
Fund waiver; I can provide any additional documentation as required.

\bigskip

\textbf{Conflicts of interest.}
None.

\bigskip

I would be grateful to have this Research Note considered for
publication, and I am available for any minor revision requested by
the editors.

\bigskip

Sincerely,

\bigskip

\noindent
Pierre-Olivier Despr\'{e}s Asselin \\
Independent Researcher \\
ORCID: 0009-0009-7611-4550

\end{document}
